% Exercise 4.1-2

\begin{algorithmic}[1]
  \TITLE{\textsc{Brute-Force-Maximum-Subarray}$(A, \mathit{low}, \mathit{high})$}
  \STATE $\mathit{maxLow} = \mathit{low}$
  \STATE $\mathit{maxHigh} = \mathit{low}$
  \STATE $\mathit{maxSum} = A[\mathit{low}]$
  \FOR{$i = \mathit{low}$ \TO $\mathit{high}$}
  \STATE $\mathit{sum} = 0$
  \FOR{$j = i$ \TO $\mathit{high}$}
  \STATE $\mathit{sum} = \mathit{sum} + A[j]$
  \IF{$\mathit{sum} > \mathit{maxSum}$}
  \STATE $\mathit{maxSum} = \mathit{sum}$
  \STATE $\mathit{maxLow} = i$
  \STATE $\mathit{maxHigh} = j$
  \ENDIF
  \ENDFOR
  \ENDFOR

  \RETURN $(\mathit{maxLow}, \mathit{maxHigh}, \mathit{maxSum})$
\end{algorithmic}

The running sum is reset at the start of each iteration of the outer
\textbf{for} loop of lines 4--11, so that $\mathit{sum}$ always holds the value
of the subarray $A[i \ltwodots j]$. Since the outer loop runs $n$ times and the
inner \textbf{for} loop of lines 6--10 runs at most $n$ times, the running time
is $\Theta(n^2)$.

As in exercise 4.1-5, $\mathit{maxSum}$ is initialised to $A[\mathit{low}]$
rather than to 0, so that the algorithm returns the largest element of an array
whose entries are all negative.
